% KDD Figure Caption for Arbitrage Frontier Visualization

\begin{figure}[t]
\centering
\includegraphics[width=\columnwidth]{figures/kdd_rational_boundary.png}
\caption{\textbf{The Arbitrage Frontier: Router Decision Boundary in Quality-Cost Space.} 
Each point represents a test prompt, colored by the router's routing decision using the production LinUCB algorithm. The dashed line represents the economic indifference curve with slope $1/\lambda = 50$, where $\lambda = 0.02$ is the cost-quality trade-off parameter for the \textit{Auto} profile. 
%
Red points (routed to Gemini-3-Pro-Preview) cluster in regions where the quality gain $\Delta Q$ justifies the cost premium $\Delta C \approx \$7/1M$ tokens. 
%
Blue points (routed to GPT-OSS-120B) represent prompts where the marginal quality gain is insufficient to justify the cost. 
%
The horizontal clustering at $\Delta C \approx \$7$ reflects the fixed pricing differential between models ($\$7.00$ vs. $\$0.06$ per 1M tokens), revealing that \textbf{the router's intelligence lies in the X-axis}: it has learned to predict that the same \$7 investment yields wildly different returns depending on prompt complexity. 
%
Models were selected from the Pareto frontier based on the Frontier Capability Index (FCI), a composite metric of HLE, GPQA, and LiveBench benchmarks, ensuring the comparison spans the full quality-cost spectrum.
%
For simple queries ($\Delta Q < 0.15$, left), the marginal quality gain is negligible and the cheap model is selected. 
%
For complex tasks ($\Delta Q > 0.15$, right), the expensive model provides substantial value and is worth the 117$\times$ cost premium. 
%
Minor deviations from the theoretical boundary result from the exploration bonus ($\alpha \times \text{uncertainty}$) inherent in LinUCB, demonstrating the algorithm's adaptive learning behavior.
%
The router achieves significant cost savings by routing most prompts to the efficient model while reserving the premium model for high-value tasks.
}
\label{fig:arbitrage_frontier}
\end{figure}

% Alternative shorter version for space-constrained venues:
% 
% \caption{\textbf{The Arbitrage Frontier.} 
% Router decisions in quality-cost space using LinUCB. 
% The dashed line shows the indifference curve (slope = $1/\lambda = 50$).
% Points cluster horizontally at $\Delta C \approx \$7$ (fixed pricing differential between
% Gemini-3-Pro-Preview and GPT-OSS-120B, selected from the FCI-based Pareto frontier), 
% while quality gains vary by prompt complexity.
% Red: routed to expensive model; Blue: routed to cheap model (117$\times$ cost difference).
% Minor boundary deviations result from UCB exploration bonus.
% }

